\cotefile{83}
\scanpage{63}
% loc: 1, top
de ce qui \uline{renverse} l'orientation. \deleted{\unsure{Ceci|Ce}}

% loc: 2–5
Donc la structure $\unsure{G'|\mathcal G'|E_i}$, dans le cas de l'icos. $p=5$, où $\deleted{\unsure{\mathbb F}}\mathbb F_5^{\times2}=\{\pm1\}$, est la structure de carte \inserted{icosaédrique} non orientée, et non celle du biicosaèdre)

% loc: left, oblique NB
% slanted
\margin{\uline{NB} l'homothétie $-1$ mérite le nom d'antipodisme, \unsure{de toutes façons}, il \deleted{\illegible{1}} \unsure{opère sur} $G'$, ssi $-1\in k^*$, -- et \unsure{existe sur|en} \unsure{autres cas}, \uline{\unsure{non}} \uline{\unsure{orienté|orientée|orientés}} \unsure{du moins de} structure $G'$. Il faudrait trouver \unsure{son|sur} expression \unsure{dans les|des} trois symétries fondamentales (dans les relations \unsure{d'un repère|du repère}$\ldots$ \unsure{suivant que} $-1\in k^*$ ou non)}

% loc: 5–8, right and centre
Ceci dit, il y a \unsure{une sol\textsuperscript{n}|le groupe} pour $G'$ par $k^\times$ opérant sur $C^*$ c'est le groupe des autom de \deleted{\unsure{cette}} \uline{structure} de multicartes.

% loc: 9–15
On trouve une représentation \unsure{naturelle|surtout|exceptionnelle} canonique remarquable des cartes du type $G_p$. \deleted{pour les} Considérons la carte universelle (``définie sur $\mathbb Z$'') qui soit triangulée ($\rho_f$ d'ordre \unsure{$3$}, \inserted{$\det\rho_f=1$, $\operatorname{Trace}\rho_f=-1$}) et pour laquelle $\rho_s$ soit \uline{unipotent} ($\det\rho_s=1$, $\operatorname{Tr}\rho_s=2$).

% loc: 16–17
dans Alg d'Azumaya engendrée \deleted{par} $\rho_a=\sigma$ d'ordre $4$, $\rho_f$ d'ordre $6$

% loc: 18–19
\[
\left\{\begin{aligned}
\det.\,\sigma&=1,&\det\rho_f&=1\\
\operatorname{Tr}\sigma&=0,&\operatorname{Tr}\rho_f&=-1
\end{aligned}\right.
\]

% loc: 20–21, bottom
avec
\[
\begin{aligned}
\sigma.\rho_f&=\deleted{\unsure{\rho_a}}\rho_s^{-1} & \rho_s\ \text{unipotent i.e.}\\
\big\vert\ \det\rho_s&=1,\ \operatorname{Tr}\rho_s=2
\end{aligned}
\]
% archive sheet number: 62

\scanpage{64}
% loc: 1–4, top
On trouve une représentation fidèle de $\mathrm{SL}(2,\mathbb Z)/{\pm1}$ \deleted{dans} (engendré par $\sigma$ et $\rho_f$ mod Centre de l'Alg. d'Azumaya.) \deleted{\unsure{du gr. dans la place|Supp. dans le plus}}. Considérons

% loc: 5–8
\deleted{Soit} la représentation de la carte triangulaire universelle, dans le plan \inserted{affine} \inserted{défini \unsure{par}} $\operatorname{Tr}u=\deleted{\unsure{2}}1$, engendrée par le \unsure{sommet} $\deleted{\unsure{\rho_s,\rho_a}}\rho_s$. \deleted{En c\unsure{ar.}}

% loc: left, diagonal at 7–14
% slanted
\margin{les sommets sont \uline{les} \uline{unipotents} les arêtes sont les paires d'unip. \unsure{dits conjugués|dits équivalents} : $(\rho_s,\underset{-\sigma\rho_s\sigma}{\underbrace{\overset{\rho_{s'},\,\text{où }s'=\sigma(s)}{\overbrace{\sigma\rho_s\sigma^{-1}}}}})$. i.e. paires d'unip\deleted{\unsure{ot}}otents $(u,u')$ tels que $\operatorname{Tr}uu'=-1$}

% loc: 9–16
En réduisant mod $n$, je dis que ça se factorise par $\mathrm{SL}(2,\mathbb Z/n\mathbb Z)/{\pm1}$ (tautologique, tel que j'ai présenté les choses), et la représentation obtenue de $C_n$ est celle que j'ai décrit tantôt (mais qui doit garder un sens, sans supposer $n$ premier $\ldots$)
% archive sheet number: 63

\scanpage{65}
% loc: 1–9, top
Soit $A$ une algèbre d'Azumaya de rang $4$ sur anneau $k$, et supposons que $\deleted{\unsure{A}}M$ soit ``triviale'', plus précisément, qu'elle provienne d'un $k$-module projectif de rang $2$ $E$, tel que le module inversible $\det E$ soit \unsure{trivial|triviale} (i.e. le \deleted{groupe structural} $\unsure{G_p(1)|\mathrm{GP}(1)|Gp(1)}_{\deleted{\unsure{k}}\,k}$-torseur décrivant $M$ peut se remonter en un $\mathrm{SL}(2)_k$-torseur.) Mais pour le moment, supposons seulem\textsuperscript{t} que $M$ soit triviale i.e.

% loc: 10
\[
 M\simeq\operatorname{End}(E).
\]

% loc: 11–14
Soit $u\in M$ un élément non dégénéré de carré nul. Alors $ku$ définit une ``droite'' $L\subset M$, et $D=ku\simeq L\otimes L'$, où $L'$ est l'orthogonal de $L$ dans $E^\vee\simeq E\otimes\deleted{\unsure{\det E}}\underline{\omega}^{-1}$ (où $\underline{\omega}=\omega(E)$).

% loc: 15–17
\deleted{Donc} En fait, sous l'identification précédente (dépendant d'une convention de signe) on aura $L'=L\otimes\underline{\omega}^{-1}$ donc

% loc: 18
\[
 D_L\simeq L^{\otimes2}\otimes\underline{\omega}^{-1},\quad
 \text{i.e.}\quad\underline{\omega}\xrightarrow[\sim]{\varphi_L}
 \underline{\operatorname{Hom}}(L^{\otimes2},D_L)
\]

% loc: 19–21
(isom can, \deleted{\unsure{à}} une fois faite la conv. de signe). Le choix d'une base $e\in\underline{\omega}$, d'où $e^{-1}\in\underline{\omega}^{-1}$ définit un iso. :

% loc: 22
\[
 \varphi_L^e:\quad\deleted{\unsure{D_L\longrightarrow{}}}\ L^{\otimes2}\xrightarrow{\sim} D_L
\]

% loc: 23–24
Si on remplace $e$ par $\lambda e$, $\unsure{e^{-1}|x^{-1}}$ est remplacé par $\lambda^{-1}e^{-1}$ et on trouve

% loc: 25
\[
 \varphi_L^{\lambda e}=\lambda\varphi_L^e
\]

% loc: 26, bottom
Quand $e$ est choisi (str. symplectique sur $E$)
% archive sheet number: 64

\scanpage{66}
% loc: 1–8, top
d'où l'iso $\varphi^e$, et si on suppose $L$ trivial i.e. $L^*$ non vide, \deleted{on a} (\unsure{donc|bien|idem} itou pour $D$, indépendamment \deleted{de} de l'hyp. qui l'avait fixée comme un $ku$) alors dans $D^*$ il y a une classe d'éléments modulo $k^{*2}$, privilégiée : celle formée par les $\varphi^e(x\otimes x)$, pour $x\in L^*$. Appelons-la $D^*(e)$ : \deleted{On aura $D^*(\lambda e)=\lambda D^*(e)$.}

% loc: 9
\[
 D^*(e)=\{\varphi_L^e(x\otimes x)\mid x\in L^*\}\qquad \text{pour }D=D_L
\]

% loc: 10–11
\deleted{donc}
\[
 D^*(\lambda e)=\lambda D^*(e)\qquad
 \left\{\begin{aligned}
 &\text{et bien sûr }\lambda^2D^*(e)=D^*(e)\\
 &\text{puisque }D^*(e)=k^{\times2}.u\ \text{pour}\\
 &u\in D^*(e)
 \end{aligned}\right.
\]

% loc: 12–14
\deleted{La subdivision de $D^*$ en des \unsure{sous}} On peut dire que pr le $k^\times$-torseur $D^*$, on a un iso can.
\[
 D^*/k^{\times2}\simeq\underline{\omega}^*/k^{*2}
\]

% loc: 15–20
\deleted{\illegible{1}} Supposons que $ {}_2\operatorname{Pic}(k)=0$, donc pour \unsure{tout|toute} \deleted{droites $L$ de} $u$ dans le cône $C^*$ des éléments de $M$ non dégénérés de carré nul, (i.e tels que $\operatorname{Tr}u=\det u=0$) la droite $L_u$ correspondante est non vide. Alors On définit \deleted{$C^*$}
\[
 C^*(e)=\bigcup_D D^*(e)
\]

% loc: 21–23, bottom
et on a ceci :
\[
 \left\{\begin{aligned}
 C^*(\lambda e)&=\lambda C^*(e)\\
 \lambda^2 C^*(e)&=C^{\unsure{2|*}}(e)
 \end{aligned}\right.\quad
 \begin{aligned}
 &\text{donc }C^*(e)\ \text{ne dépend \unsure{que|pas}}\\
 &\text{de }e\bmod k^{\times2}\deleted{\unsure{k^*}}
 \end{aligned}
\]
% archive sheet number: 65

\scanpage{67}
% loc: 1–3, top
et on trouve que
\[
 C^*=\coprod_{e\in\underbrace{\underline{\omega}^*/k^{*2}}_{\substack{\text{torseur}\\\text{sous }k^\times/k^{\times2}}}}C^*(e)
\]

% loc: 4–11
partition qui est permutée transitivement par les homothéties \inserted{de $M$} : \deleted{l'on \unsure{que|pour}} $\lambda C^*(e)=C^*(\lambda e)$. L'ens des classes de cette partition \deleted{est} \deleted{\unsure{prend donc un}} de façon naturelle (par la relation d'homothétie) un torseur sous $\unsure{k^{\times2}/k^{\times2}|k^\times/k^{\times2}}$, et ce torseur est can. iso à $\underline{\omega}^\times/k^{\times2}$.

% loc: 12–16
Cette sous-division de $C^\times$ en classes ne dépend que de $M$, et pas du choix de $E$ (avec $\det E\simeq k$) et d'un iso $\operatorname{End}(E)\simeq M$. On peut aller plus loin, \deleted{\unsure{quand on On}}.

% loc: left at 12–17, with arrow to Non
% slanted
\margin{On doit supposer ici $ {}_2\operatorname{Pic}(k)=0$ ; (Que se passe-t-il autrement ?) Il semble que \unsure{même en|Mais on a} $=0$ \unsure{de|pour} cette subdivision de $C^*$ est invariante par le groupe \unsure{privé|propre} $\operatorname{Aut}(M)$ \unsure{des|de}. $\longrightarrow$ \uuline{Non}}

% loc: 17–19
On \deleted{\unsure{aura}} a une application surjective
\[
 \begin{aligned}\varphi^e:\ E^*&\longrightarrow C^*(e)\\
 x&\longmapsto\varphi^e(x\otimes x)\end{aligned}
\]

% loc: 20–22, bottom
satisfaisant \unsure{à}
\[
 \left\{\begin{aligned}
 \varphi^e(\lambda x)&=\lambda^2\varphi^e(x)\\
 \varphi^{\lambda e}(x)&=\lambda\varphi^e(x)
 \end{aligned}\right.
\]
% archive sheet number: 66

\scanpage{68}
% loc: 1–5, top, including interlinear additions
Ainsi, \deleted{\unsure{les}} \inserted{chaque} classe $C_i^*\deleted{\unsure{(e)}}$ \inserted{$(=C^\times(e))$} est munie \deleted{\unsure{chacune}} \deleted{(pour $e$ fixé)} d'une structure \unsure{triangulée|triangul.|triangulaire} \deleted{\unsure{non}} orientée canonique, \inserted{(indépendante du choix de $e$)} avec en plus une $\lambda$-orientat. \inserted{(dépendant à priori du choix de $e$)} qu'il faudrait expliciter. Le groupe projectif de $M$

% loc: 6
\[
 M^*/k^\times\quad(\simeq\operatorname{End}(E)/k^\times)
\]

% loc: left, oblique note beside 6–12
% slanted
\margin{c'est une pure \unsure{pétit\textsuperscript{o}|pét.|preuve}, il faut voir que si $u,v\in C^*$ sont \deleted{\unsure{tels}} tels que $\operatorname{Tr}uv=-1$, ils appartiennent à une $=$ classe $\ldots$ \unsure{cf|c'est} plus bas $\ldots$}

% loc: 7–8
permute entre \deleted{\unsure{elles|deux}} eux les $C_i^*\deleted{\unsure{(e)}}$, en respectant \inserted{(par transport des structures)} les structures orientées.

% loc: 9–13
On fera attention que pour \deleted{les autom qui \unsure{proviennent de}} $u\in\mathrm{Gl}(E)$, définissant $\bar u\in\mathrm{Gl}(E)/k^\times\simeq M^\times/k^\times$, l'endom. de $\deleted{\unsure{\mathrm{Gl}^{\times}(E)|\mathrm{Gl}(E)}}M^*$ \deleted{\unsure{qu}} (donc de $C^\times$) qu'il définit, est donné par
\[
 \theta\longmapsto u(\theta)=u\theta u^{-1}
\]

% loc: left at 14–18, double vertical rule
% slanted
\margin{\uline{Attention}}

% loc: 14–18
Ce n'est compatible avec l'application
\[
 \varphi^e:E^*\xrightarrow{\sim} C^\times(e)\subset C^*
\]
(\unsure{avec|car} $\mathrm{Gl}(E)$ opérant sur $E^\times$ et sur $C^\times$) que sur le sous-groupe $\mathrm{Sl}(E)$. De façon précise, on aura

% loc: 19, boxed
\[
 \boxed{\theta(\varphi^e(x))=\det\theta\,(\varphi^e(\theta(x)))}
\]

% loc: 20–21, bottom
\[
 \bigl(\ \theta(x\overset{e}{\otimes}x)=
 \theta(x)\overset{\theta(e)}{\otimes}\theta(x)=
 \theta(x)\overset{\det\theta.e}{\unsure{\odot|\otimes}}\theta(x)
 =\det\theta\bigl(\theta(x)\overset{e}{\otimes}\theta(x)\bigr)
\]
\deleted{Donc le groupe}
% archive sheet number: 67

\scanpage{69}
% loc: 1–6, top
\uline{Remarque} Indépendamment de tte hypo sur $\operatorname{Pic}(k)$ et sur $M$, on peut regarder $C^*$, et la structure de graphe sur $C^\times$ définie par la relation $\operatorname{Tr}(uv)=-1$, d'où structure triangulée. Mais

% loc: 7–9
1\textsuperscript{o}) $C^*$ \deleted{\unsure{ne|n'est}} n'est non vide que si $\deleted{\unsure{k}}M$ est triviale, plus précisément si $M\simeq\operatorname{End}(E)$ avec $E=k\oplus L$ ($L$ module inv.)

% loc: 10–17
\deleted{On a dit} 2\textsuperscript{o}) \deleted{Il n'a \unsure{Il|Je}} L'ens des arêtes est non vide ssi $M\simeq M_2(k)$. En effet, si $\exists$ arête non vide $\{u,v\}$ \deleted{$\operatorname{Tr}uv=$} avec $u,v\in C^\times$, $\operatorname{Tr}uv=-1$, les théorèmes de conjugaison nous fournissent un iso unique $\mathbb M_2\longrightarrow M$, transformant $u_0$ en $u$, $v_0$ en $v$ ($(u_0,v_0)$ couple type dans $M$, satisfaisant les $\unsure{={}^s|\text{mêmes}|=}$ conditions \unsure{:|.}).

% loc: 18–20, bottom
Donc on a intérêt à se borner \unsure{alors|tjs} à des iso $M\simeq\operatorname{End}(E)$, en supposant $E$ \uline{libre}
% archive sheet number: 68

\scanpage{70}
% loc: 1–2, diagonally struck
\deleted{utilisant la formule}
\[
 \deleted{(a\overset{e}{\otimes}\underbrace{b)\circ(c}\overset{e}{\otimes}d)
 \simeq c\overset{e}{\otimes}b\,(d\wedge a/e),}
\]

% loc: 3–5, same cancellation
\deleted{on trouve}
\[
\begin{aligned}
 &\deleted{uvw=(y\overset{e}{\otimes}x\,\underbrace{(-x\wedge y/e)}_{-1})\,z\overset{e}{\otimes}z}\\
 &\deleted{\phantom{uvw}=-z\otimes x\,\underbrace{(z\wedge y/e)}_{-1}=z\overset{e}{\otimes}x}
\end{aligned}
\]

% loc: 6–8
Il existe donc un unique $\rho\in\mathrm{Gl}(E)$ tel que
\[
 \rho x=y,\ \rho y=z,\ \rho z=x
\]

% loc: 9–13
et ce $\rho$ est d'ordre 3. Il ne change pas en remplaçant $\unsure{\{x,y,z\}|(x,y,z)}$ par $(-x,-y,-z)$, il ne dépend que du triplet $\{\dot{x},\dot{y},\dot{z}\}$ et d'une orientation de ce triplet. Changer l'orientation, revient à changer $\rho$ en $\rho^{-1}$

% loc: 14–15
\unsure{Or on aura|On en aura|On a aussi}
\[
 \det\rho=1,\ \operatorname{Tr}\rho=1,\ \inserted{\text{donc}}\quad
 \rho^{-1}=\deleted{\unsure{\rho}}1-\rho=\rho'\quad(\text{conj. de }\rho)
\]

% loc: 16–17
Considérons
\[
 \unsure{\sigma_f|\sigma|\tau} =x\overset{e}{\otimes}x+y\overset{e}{\otimes}y+z\overset{e}{\otimes}z
 \qquad\text{On a }\left|\begin{aligned}
 \operatorname{Tr}\sigma&=0\\
 \det\sigma&=\deleted{4}3.
 \end{aligned}\right.\quad
 \text{(par formule d'additivité du dét\ldots)}
\]

% loc: 18–20
c'est un endom qui commute à $\rho$, donc de la forme
\[
 \sigma=a\rho+b
\]

% loc: 21–24, bottom
\deleted{\unsure{où}} d'où, en prenant les traces
\[
\begin{aligned}
 &\operatorname{Tr}\sigma.=a+2b\qquad\text{donc }\underset{\deleted{\unsure{b=-a}}}{a=-2b},\ \text{donc}\\
 &\sigma=\deleted{\unsure{a(-2b|a\rho+2b}}\,b(-2\rho+1)\\
 &\det\sigma=b^2\,\underline{\det(1-2\rho)},\qquad \text{or }\det(1-2\rho)=3
\end{aligned}
\]
% archive sheet number: 69

\scanpage{71}
% loc: 1–2
Donc $b^2=1$ i.e. $b=\pm1$ donc
\[
 \sigma=\pm(1-2\rho)
\]

% loc: 3–4
\deleted{Supposons qu'on ait}
\[
 \sigma=1-2\rho
\]

% loc: 5–8
alors on ne peut avoir en $=$ temps $-(1-2\rho)$ (car $2\,1_k\ne0$), donc cela fixe $\rho$ dans la \deleted{paire} $\{\rho,\rho'=1-\rho\}$. Mais il faut voir lequel ! \deleted{Donc quitte à rectifier le tir, écrivons}

% loc: 9, crossed formula
\[
 \deleted{u+v+w=1-2\rho\qquad \rho\in M^*,\ \det(\rho)=1,\ \rho^3=1}
\]

% loc: 10–15
Revenant : $M$, on peut dire que pour un triangle $\{u,v,w\}$ de $C^*(e)$, existe $\rho\in M^*$ avec
\[
 \left\{\begin{aligned}
 \det\rho&=1&\rho&\ \text{induit une permutation}\\
 &&&\text{circulaire entre les }u,v,w\\
 \rho^3&=1&&\text{et on aura }\operatorname{Tr}\rho=\operatorname{Tr}\rho'=1
 \end{aligned}\right.
\]

% loc: left of conditions, oblique
% slanted
\margin{C'est la relation $\operatorname{Tr}\rho=1$ qui est la plus simple pour lever l'ambiguité mod scalaire (une fois que la perm. \unsure{ass.} est fixée), \deleted{\unsure{dét}} \unsure{les rel.} $\det\rho=1$, $\rho^3=1$ viennent par surcroît $\ldots$}

% loc: 16–18
\unsure{que|qui} $\rho$, à passage à $\rho'$ près. Il y a cependant un seul tel $\rho$, tel que
\[
 \boxed{u+v+w=1-2\rho}
\]
\margin{vérifier si c'est le bon !}

% loc: 19–22, bottom
D'où la structure \deleted{\unsure{de}} orientée : \inserted{pour} $\{u,v,w\}$ formant triangle, \deleted{\unsure{le}} le $\rho$ qu'on vient de décrire définit sur l'ens $\{u,v,w\}$ une permutation circulaire canonique.
% archive sheet number: 70

\scanpage{72}
% loc: 1
et on a un \deleted{bijection} diagramme commutatif

% loc: 2–5, upper diagram
\diagram{\begin{tikzcd}[ampersand replacement=\&,column sep=large,row sep=huge]
 {\left[E^*/k^*\underset{\substack{\vert\\\text{si }\operatorname{Pic}(k)=0\\\text{sinon }\hookrightarrow}}{\simeq}\mathbb P(E)(k)\right]}
 \arrow[r,"{\substack{\text{induit}\\\deleted{\text{\unsure{fonctoriel|fonctif}}}\,\text{iso de}\\\text{schémas}\,\deleted{\unsure{\sim}}\\\sim}}","{\text{bien connu}}"']
 \& C^*/k^* \& {}\\
 E^*\arrow[u]\arrow[r,"{\varphi^e}"]
 \& C^*(e)\arrow[u]\arrow[r,phantom,"{\subset}"]
 \& C^\times\arrow[ul,dashed]
\end{tikzcd}}

% loc: 6–7
On a $\varphi^e(x)=\varphi^e(y)$ \unsure{ses|sés|ssi} $y=\lambda x$ avec $\lambda\in\mu_2(k)$ (i.e. $\lambda\in k^\times$, $\lambda^2=1$).

% loc: 8–9
Supposons maintenant
\[
 \{\pm1\}\xrightarrow{\overset{\text{iso}}{\sim}}\mu_2(k).
\]

% loc: 10–16
[Si 2 inv. Cela signifie que $\operatorname{Spec}k$ est 0-connexe. Si $\operatorname{Spec}k$ a un nb fini \inserted{$n$} de comp. conn. \deleted{\unsure{soient}}, cela signifie que \deleted{toutes} le nb de comp. connexes où $2$ est nul est $n-1$ -- i.e.$^*$ $k=k'\times k''$ \deleted{\unsure{avec}} \unsure{anneau produit|anneau premier}, avec $2$ nul dans $k'$, \uline{non} nul dans $k''$, et $\operatorname{Spec}(k'')$ est \deleted{\unsure{co}} connexe.]

% loc: left beside 13–16
% slanted
\margin{le i.e. est sans doute vrai sans hyp. sur le nb de comp. connex.}

% loc: 17–20, bottom
Alors l'hom induit
\[
 \psi^e:\quad E^*/{\pm1}\xrightarrow{\sim}C^*(e)
\]
est \uline{bijectif}. Cela nous permet de transporter sur $C^*(e)$ la structure
% archive sheet number: 71

\scanpage{73}
% loc: 1–3
combinatoire $I_e$ sur $E^*/{\pm1}$ : structure de \deleted{graphe} \unsure{$k$|$\Delta$|bi}-polyèdre \uline{orienté}. Je rappelle que cette structure consiste en

% loc: 4–5
a) Structure de graphe : $\{\dot x,\dot y\}$ forment une arête \deleted{\unsure{si}} sss $\overbrace{x\wedge y}=e$ ou $-e$ ;

% loc: 6–8
$\longmapsto$ structure triangulée $\{\unsure{x|\dot x},y,z\}$ est un triangle ssi $\{\dot x,\dot y\}$, $\{\dot y,\dot z\}$, $\{\dot z,\dot x\}$ sont des arêtes.

% loc: 8–9, NB between a and b
\inserted{\uline{NB} a) Chaque arête est contenue dans exactement $\underline{2}$ triangles b) les triangles sont les triplets de la forme $\{\dot x,\dot y,\dot z\}$, avec $x+y+z=0$,}

% loc: 9–13
b) Chaque triangle est \uline{orienté} -- \deleted{(l'arête $\{\dot x,\dot y\}$ est directe (N} l'arête $\{\dot x,\dot y\}$ du triangle défini par $x,y,z$ avec $x\unsure{-|+}y+z=0$, \deleted{et $x\wedge y$} est \uline{directe}, sss $x\wedge y=e$, d'où $y\wedge z=e$, $z\wedge x=e$

% loc: 13–14, interlinear
\inserted{De telle façon que les deux or. d'une arête induites par celles des triangles adjacents soient opposées}

% loc: 15–24, lower part
c) L'``ens des arêtes et triangles incidents à une arête, qui est un contour combinatoire (\unsure{pas|non|un} néc. connexe) sur lequel $\mathbb Z$ opère (de par les données structurelles précédentes) est muni'' de plus d'une opération de $k$ dessus, prolongeant celle de $\mathbb Z$, et pour laquelle l'ens des arêtes \inserted{($\longmapsto$ sommets)} et l'ens des \deleted{triangles} ($\longmapsto$ arêtes) sont des \inserted{$\unsure{k|\lambda|\mathbb Z}$-torseurs}.
% archive sheet number: 72

\scanpage{74}
% loc: 1–4
\uline{NB} Si $s=\dot x$, \deleted{et si $a$} on définit $T_\lambda=T_\lambda^s$ sur l'ens des $t$ liés à $s$, i.e. des $\dot y$ avec $x\wedge y=e$, par
\[
 T_\lambda(y)=\dot{\overline{y+\lambda x}}
\]

% loc: 5, cancelled
\deleted{(d'où \unsure{$\dot y$} $T_\lambda(y)=(\unsure{-}$}

% loc: 6–9
(Si on remplace \unsure{$a$|$x$} par $-x$, on trouve la $=$ opération $T_\lambda$) En fait, ici \deleted{T} l'opération $T_\lambda$ est induite par \unsure{la|les} transvection de $E$ : \deleted{\unsure{$cx$|$E$} $u\longmapsto{}$} $\mathrm{id}+\varphi^e(x\otimes x)$

% loc: 10–11
Cette structure se transporte sur $C^*(e)$. Il vaut la peine de l'expliciter :

% loc: 12–13
Notons que
\[
 \underbrace{x\overset{e}{\otimes}x}_{=\varphi^e(\deleted{\unsure{x}}x)}
 \circ\underbrace{y\overset{e}{\otimes}y}_{=\varphi^e(y)}
 =y\overset{e}{\otimes}x\,((\unsure{-x|x}\wedge y)/e)
\]

% loc: right of composition
\margin{\uline{NB} $\operatorname{Tr}\,\deleted{\unsure{x}}y\overset{e}{\otimes}x=(\unsure{x|a}\wedge y)/e$}

% loc: 14–15
donc
\[
 \operatorname{Tr}(\deleted{\unsure{x\otimes}}\underset{\deleted{\unsure{x\overset{e}{\otimes}x}}}{u}\circ v)
 =-\bigl((x\wedge y)/e\bigr)^2
\]

% loc: left at 14–19
% slanted
\margin{Que se passerait-il si on choisit sur $C^*$ la structure de graphe définie par ``l'autre'' relation ``non dégénérée pour ttes caractéristiques'' $\operatorname{Tr}uv=+1$ ??}

% loc: 16
donc
\[
 x\wedge y=\pm e\quad\Longleftrightarrow\quad\boxed{\operatorname{Tr}uv=-1}
\]

% loc: 17–19
Les triangles $\{u,v,w\}$ de $C^\times(e)$ sont donc les triples $\{u,v,w\}$ tels que
\[
 \operatorname{Tr}uv=\operatorname{Tr}vw=\operatorname{Tr}wu=-1
\]

% loc: 20–21, bottom edge
Si \deleted{\unsure{on}}
\[
 \begin{aligned}
 x+y+z&=0 & x\wedge y&=y\wedge z=z\wedge x=e\\
 u&=x\overset{e}{\otimes}x,\quad v=y\overset{e}{\otimes}y,\quad & w&=z\overset{e}{\otimes}z=e
 \end{aligned}
\]
% archive sheet number: 73

\scanpage{75}
% loc: 1, upper right
Azumaya rang 4

% loc: 2, below title
avril 1986
% archive sheet number: 74

\scanpage{76}
% loc: 1–4, top, with polynomial substitutions
\[
\begin{aligned}
 \delta(\underline{xyzt})&=\delta(xy(zt))=Q\bigl(\delta(x),\delta(y),\delta(zt),\delta(xy),\underline{\delta(yzt)},\underline{\delta(ztx)}\bigr)\\
 &\deleted{\ne\delta(x(yz),t)}\\
 &=\delta(x(yz)t)=Q\bigl(\delta(x),\delta(yz),\delta(t),\underline{\delta(xyz)},\underline{\delta(yzt)},\delta(tx)\bigr)\\
 &=\delta((xy)zt)=Q\bigl(\delta(xy),\delta(z),\delta(t),\underline{\delta(xyz)},\delta(zt),\underline{\delta(txy)}\bigr)
\end{aligned}
\]
\margin{\unsure{remplacer par|rempl. par|remplaçons les} Pol. en $\delta(y),\delta(z)\,\unsure{\&|,}\,\delta(t)$}

% loc: 5–7, beneath equations
\diagram{\begin{tikzcd}[ampersand replacement=\&,column sep=small,row sep=small]
 {} \& yzt\arrow[dr,dash] \& \deleted{\unsure{z}}ztx\\
 xyz\arrow[dr,dash] \& {} \& yzt\\
 {} \& xyz \& txy
\end{tikzcd}}

% loc: 8–11
On trouve trois expressions polynômiales différentes d'une $=$ quantité, d'où 2 relations (\unsure{prob\textsuperscript{t}|probl\textsuperscript{t}|quasi !} \unsure{équivalentes|arguables}) entre les dix quantités
\[
 \delta(x),\delta(y),\delta(z),\delta(t),\quad\delta(xy)\ \delta(xz)\ \delta(xt)\ \delta(yz)\ \delta(yt)\ \delta(zt)
\]

% loc: 12–14
Exprimer que si $\tau$ normalise $H\subset G$, $H$ \uline{comm.}, alors $h\in H$,
\[
 \boxed{\delta(\tau h)\overset{!}{=}\delta(\tau)}
\]
i.e. les $\tau^n(h)\inserted{\overset{\mathrm{déf}}{=}\tau^nh\tau^{-n}}$ commutent à $h$ ($n\in\mathbb Z$), alors on a la formule encadrée. On aimerait l'écrire comme une ``identité'' inconditionnelle

% loc: left, oblique beside 12–17
% slanted
\margin{C'est \uline{idiot}, \unsure{comme on voit|comme on sait} en prenant $\tau\in H$ -- il faut supposer de plus $\tau$ \unsure{ou} ``$\tau$ moins \unsure{un}'' opère de façon \unsure{non triviale|trivialement} sur $H$, i.e. qque chose du goût $\ldots$}

% loc: 15–17
\[
 \delta(\tau h)-\delta(\tau)\overset{?}{=}
 \text{expression polynomiale}\ \inserted{(\text{à coeff}\in\mathbb Z)\ \deleted{\text{\unsure{système constant|système constante}}}}\ \text{en les}
\]
\[
 \delta\bigl(\underbrace{\{\tau^n(h),h\}}_{\text{commutateur}}\bigr)
 \ \text{et }\delta(\tau),\delta(\tau h),\delta(h)
\]
\inserted{dont tous les termes contiennent un $\delta(\{\cdot,\cdot\})$}

% loc: 18–19
(peut-être suffit-il de
\[
 \underbrace{\{\tau h\tau^{-1},h\}}_{\tau h\tau^{-1}h\tau h^{-1}\tau^{-1}h^{-1}}
 \ \text{et}\quad
 \underbrace{\{\tau^{-1}h\tau,h\}}_{\tau^{-1}h\tau h\tau^{-1}h^{-1}\tau h^{-1}}
\]

% loc: 20, bottom edge
A-t-on au moins
\[
 \boxed{\delta(\tau.\underbrace{\tau.(h)h^{-1}}_{\{\tau,h\}})\overset{?}{=}\delta(\tau)}
 \quad(\text{si }\tau\ \text{normalise }H\ \text{comm, }h\in H)
\]
% archive sheet number: 75

\scanpage{77}
% loc: 1, upper centre
\mspage{PU 93}

% loc: 2, upper left
\deleted{$\mathrm{IV}^{\prime\prime}_{4}$\quad E v o l u t i o n}

% loc: right margin, opposite first list
\margin{\handwritten{\unsure{,|)}}}

% loc: 3–5
\deleted{E s s o r - d é c l i n}\\
\deleted{croissance - vieillissement}\\
\deleted{régénération - usure}

% loc: 6–8
\deleted{enfance (ou jeunesse) - vieillesse}\\
\deleted{innocence - maturité}\\
\deleted{l ' e n f a n t - l e\quad v i e i l a r d}

% loc: 9–12
\deleted{naître - mourir}\\
\deleted{naissance - mort (*)}\\
\deleted{création - d\deleted{\unsure{e|é}}\handwritten{e\unsure{s}}truction}\\
\deleted{apprendre - \unsure{desapprendre|désapprendre} (V)}

% loc: 13–16
\deleted{commencement - fin}\\
\deleted{origine - destination}\\
\deleted{départ - retour}\\
\deleted{sortir - rentrer}

% loc: 17–21
\deleted{le tôt - le tard}\\
\deleted{le précoce - le tardif}\\
\deleted{matin - soir}\\
\deleted{princtemps - automne}\\
\deleted{est - ouest}

% loc: 22, second heading
\deleted{\deleted{I}\handwritten{\unsure{4|/}}$\mathrm{V}^{\prime\prime}_{1}$\quad H a u t e u r}

% loc: 23–27
\deleted{L e\quad h a u t - l e\quad b a s}\\
\deleted{monter - descendre}\\
\deleted{montéee - descente}\\
\deleted{élévation - profondeur (**)}\\
\deleted{l e\quad c i e l - l a\quad t e r r e}

% loc: 28–30
\deleted{hauteur (ou longueur) - largeur (**)}\\
\deleted{le vertical - l'h\deleted{\unsure{r|o}}orizontal}\\
\deleted{maigreur/ - corpulence}

% loc: 31–33
\deleted{étendue - profondeur}\\
\deleted{le vaste - le profond}

\deleted{l'aigu - le grave}

% loc: 34–35, first footnote
\rule{1.25cm}{0.3pt}
\deleted{(*) Comparer ce couple avec le couple voisin ``vie - mort'', que j'ai inclus dans le groupe II (``action - inaction'').}

% loc: 36–39, second footnote
\deleted{(**) Comme tous les couples yin-yang qui ont une double acception, l'\unsure{une|un} au sens propre et l'autre au sens figuré, les couples ``élévation - profondeur'' et ``hauteur - largeur'' peuvent s'entendre dans l'un et l'autre sens.}
% archive sheet number: 76

\scanpage{78}
% loc: 1–3, boxed upper right
\fbox{\parbox{.92\linewidth}{Description \deleted{du} \inserted{de l'anneau} $\mathcal A\simeq M_2(k)$ ($k$ corps) en termes \unsure{de|du} $G=\mathcal A^\times/k^\times$ (tentative)}}

% loc: 1–5, left
\[
\begin{aligned}
1*1&=1\\
(x\deleted{\unsure{s|a}}*y)(y*x)&=1\\
xs*y&=x*y\quad\text{si } s,y\text{ commutent}\\
[\,x*yt&=x*y\quad\text{si }t,x\quad\text{---}\,]\\
[\,x*y&=1\quad\text{si } x,y\text{ commutent}\,]
\end{aligned}
\]

% loc: 3–5, right
\[
\begin{aligned}
[\,x*y&=\bar{x}\,\bar{y}\,\bar{x}^{-1}\bar{y}^{-1}\\
\delta(\dot x)&=\deleted{\unsure{\mathrm{discriminant}|\mathrm{discrim}}}\frac{b^2-4c}{c}\,]\\
c&=\det\dot x,\ b=\operatorname{Tr}x
\end{aligned}
\]

% loc: 6–7
\[
(x*y)^2-P(\delta x,\delta y,\delta(xy))\,x*y+1=0\qquad P\in\mathbb Z[X,Y,\unsure{Z|\mathbb Z}]
\]
\margin{bien déterminé}

\[
[\delta(z),x*y]=[\delta(z),\delta(z')]=0
\]

% loc: 8–11, framed group and annotation below
\[
\boxed{\left\{\begin{aligned}
\delta(xy)&=\delta(yx)\qquad\text{i.e. }\delta(g(x))=\delta(x)\qquad(x,y,g\in G)\\
\delta(x)&=\delta(x^{-1})\\
\delta(xy)+\underbrace{\delta(xy^{-1})}_{\shortparallel\,\delta(x^{-1}y)}
 &=\deleted{\unsure{\delta(x)\delta(y)}}(\delta(x)+2)(\delta(y)+2)-4
 =\delta(x)\delta(y)+2(\delta(x)+\delta(y))
\end{aligned}\right.}
\]
\margin{si $x,y$ commutent seulement ??? À vérifier !}

% loc: 12–13
\[
[\,\boxed{\delta(1)=0}
\]
\[
\delta(x^n)=P_n(\delta(x)),\qquad\delta(\underbrace{x^ny^m}_{\text{ou d'autres monômes non commutatifs\ldots}})=P_{n,m}(\delta(x),\delta(y),\delta(xy))\quad\text{etc.}
\]

% loc: 14–18
$x_1\ x_2\ x_3\ x_4,\in G$, il doit y avoir une relation pol. à coefficients dans $\mathbb Z$ entre
\[
\delta(x_i)\quad(1\leq i\leq4),\quad\deleted{\unsure{\delta}}\quad\delta(x_ix_j)\quad(1\leq i<j\leq4)
\]
(dix quantités, \deleted{\unsure{pour en un ``espace''|pour sur un ``espace''|pour un ``exprimant''}} dépendant de ``neuf paramètres''\ldots\ ]

% loc: 19–20, boxed
\[
\boxed{\delta(xyz)=\underset{\text{pol. univ. à coeff. entiers en six variables}}{Q}(\delta(x),\delta(y),\delta(z),\delta(xy),\delta(yz),\delta(zx))\qquad ?}
\]
\margin{à déterminer}

% loc: 21–24
\deleted{Ceci donne $\delta(x^n)$ comme pol.}
\[
\begin{aligned}
\delta(x^2)\deleted{\unsure{4|+}}&=(\delta(x)+2)^2-4=\delta(x)^2+4\delta(x)\\
\delta(x^n)&=P_n(\delta(x))\qquad n\in\mathbb Z\qquad P_n=P_{-n}\\
\delta(x^ny^m)&=P_{\unsure{n,m|m,m}}(\delta(x),\delta(y),\delta(xy))
\end{aligned}
\]
\margin{[cas particulier de $\delta(xy)+\delta(xy^{-1})=\ldots$]}

% loc: 25–28
\[
\delta\bigl(\deleted{\text{monôme en}}\,\inserted{M(x_1,\ldots,x_n)}\bigr)=P_M\bigl((\delta(x_i)),(\delta(x_ix_j))\bigr)
\]
\margin{monôme non commutatif à exposants entiers $\in\mathbb Z$}
\margin{pol. à coeff. dans $\mathbb Z$ (pas uniquement déterminé si $n\geq4$)}
% archive sheet number: 77

